\section{Relation to established theories}\label{sec:theories}
Each comparison is classified as \textbf{identical} (an established equivalence), \textbf{special case} (a proved inclusion), \textbf{analogy} (similar behaviour without a proved equivalence) or \textbf{open}. Nothing in the manuscript is proved by appeal to these theories, and no claim of priority is made.
\begin{description}[leftmargin=1.2em,itemsep=3pt]
\item[Arithmetic dynamics and divisibility systems.] The map $C$ is a rational map of the positive rationals defined by arithmetic data; its occurrence set is a divisor condition on the numerator. \emph{Analogy}: we do not connect $C$ with the standard theory of dynamics of rational functions; the comparison with $p$-adic or arithmetic-dynamical invariants is \emph{open}.
\item[Geometry of numbers and lattice relations.] The lattice $\Lambda_b\subseteq\Z^k$ of exponent vectors with product $\equiv1\bmod b$ is the kernel of a homomorphism to $(\Z/b)^\times$; its transport (\Src{I}{thm:P2T1}) is a statement about kernels of homomorphisms. \emph{Identical} with the standard description of such sublattices (cf.\ Cassels, \emph{An Introduction to the Geometry of Numbers}); the dimensional results are UNNS-specific.
\item[Graver bases.] The set $\Gr_q$ is the Graver basis of the lattice restricted to a box \cite{Graver1975,Sturmfels1996}, and the persistence criterion \Src{I}{thm:P2T2} is a statement about conformal minimality. \emph{Special case} of the notion of a Graver basis; the monotone-minimality principle (\cref{lem:minimal}) is a general, elementary statement which these criteria instantiate.
\item[Rewriting systems and confluence; trace monoids.] Swap classes are the classes of a rewriting system generated by commutation of adjacent letters when both orders are admissible, as in partially commutative (trace) monoids \cite{DiekertRozenberg1995}. The commutation here is \emph{state-dependent}, so a history space is not a free partially commutative monoid; \emph{analogy}. Confluence of the block rewrites is studied in Paper~II; whether the swap graph can be described by a trace-type normal form is \emph{open}.
\item[Cubical and cellular complexes; homology.] The complexes $K(D)$ are cubical in the sense of Paper~III's definition; homology is computed over $\mathbb F_2$ (and checked over another field in the sources), and cellular homology and the long exact sequence of a pair are standard \cite{Hatcher2002}. \emph{Identical} use of standard tools; the divisor-interval cell structure is specific.
\item[Zig-zag and filtered structures.] The shape $K\leftarrow K^\circ\to K'$ of an event resembles a zig-zag diagram of complexes \cite{CarlssonDeSilva2010}; we use only that shape. \emph{Analogy}; no use is made of zig-zag persistence theory.
\item[Information-losing projections.] The factorization lemma \cref{lem:factor} is the universal property of a quotient by the fibre relation of a surjection; \emph{identical}.
\end{description}
